
\section{Pràctica de laboratori: Caçadors de forces}

Aplicarem el nostre mètode: si un objecte canvia la seva velocitat o la direcció del moviment, hi ha d'haver una força al darrere.

Avui ens convertirem en caçadors de forces. Intentarem localitzar-les, observar com actuen i elaborar un petit catàleg del seu comportament.

De moment no necessitem fórmules complicades. Observarem, farem mesures i prendrem nota. Més endavant intentarem explicar tot el que hem descobert. Una feina que ha ocupat els físics durant uns quants segles.

\subsection{1. El nostre detector de forces}

Un \textbf{dinamòmetre} és un aparell que mesura forces en newtons (N). Funciona gràcies a una molla que es deforma quan estirem d'ella.

\begin{enumerate}
    \item Observa el dinamòmetre i la seva escala.
    \item Estira suaument del ganxo. Observa el moviment de l'indicador.
    \item Penja-hi un objecte i anota la lectura.
    \item Retira l'objecte i estira amb la mà fins a obtenir la mateixa lectura.
    \item Repeteix l'experiment amb un altre objecte.
\end{enumerate}

\textbf{Anota:} Quines lectures has obtingut? Has aconseguit el mateix resultat de dues maneres diferents?

\textbf{Primera descoberta:} forces d'origen diferent poden produir la mateixa deformació. Fins i tot l'indicador del nostre detector necessita forces per moure's.

\subsection{2. Dues forces, una sorpresa}

Muntem el dinamòmetre en un bastidor i pengem una massa d'un dels extrems.

\begin{enumerate}
    \item Anota la lectura del dinamòmetre.
    \item Ara col·loca el dinamòmetre entre dues cordes que estiren en sentits oposats, mitjançant dues masses iguals i politges.
    \item Abans de mirar l'escala, escriu quina lectura esperes.
    \item Comprova el resultat.
\end{enumerate}

\textbf{Predicció:} \rule{1.5cm}{0.4pt} N

\textbf{Lectura real:} \rule{1.5cm}{0.4pt} N

Les dues masses estiren, però el dinamòmetre no marca el doble!

\textbf{Què passa?} Les forces oposades mantenen el sistema en equilibri. El dinamòmetre mesura la tensió, no la suma de les dues forces oposades.

En el primer muntatge, el bastidor també estirava del dinamòmetre. Simplement, no ens havíem fixat en ell.

\subsection{3. Un viatge costa amunt}

Col·loca un carretó amb rodes sobre un pla inclinat i subjecta'l amb el dinamòmetre, paral·lelament a la superfície.

\begin{enumerate}
    \item Comença amb el pla horitzontal.
    \item Augmenta la inclinació dues o tres vegades.
    \item Mantén el carretó immòbil i anota la força que marca el dinamòmetre en cada cas.
\end{enumerate}

\begin{center}
\begin{tabular}{|c|c|}
\hline
Inclinació & Força (N) \\ \hline
Horitzontal & \\[0.7em] \hline
Petita & \\[0.7em] \hline
Mitjana & \\[0.7em] \hline
Gran & \\[0.7em] \hline
\end{tabular}
\end{center}

\textbf{Observa:} Canvia la massa del carretó? Canvia el seu pes? Què passa amb la lectura del dinamòmetre?

\textbf{Una pista:} el pes és la força gravitatòria que exerceix la Terra sobre un objecte. La massa, en canvi, és una propietat de l'objecte. No són el mateix.

\subsection{4. Una força que no havíem convidat}

Substitueix el carretó per un bloc de fusta que fregui considerablement amb la superfície.

\begin{enumerate}
    \item Estira suaument del bloc amb el dinamòmetre.
    \item Augmenta progressivament la força fins que comenci a lliscar.
    \item Repeteix l'experiment amb diferents inclinacions.
    \item Deixa el bloc sobre el pla sense estirar-lo.
\end{enumerate}

\textbf{Observa:} Pot existir una força aplicada sense que el bloc es mogui? A partir de quina inclinació comença a lliscar?

Hem trobat el \textbf{fregament}, una força que s'oposa al lliscament o a la tendència a lliscar entre superfícies.

El fregament estàtic pot adaptar-se fins a un límit. No és una força que tingui sempre el mateix valor.

\subsection{5. La força que desapareix}

Penja un objecte d'una corda i deixa'l en repòs.

Sobre l'objecte actuen dues forces: el \textbf{pes}, cap avall, i la \textbf{tensió}, cap amunt. Com que està en equilibri, les dues forces es compensen.

\begin{enumerate}
    \item Observa l'objecte immòbil.
    \item Dibuixa dues fletxes que representin el pes i la tensió.
    \item Sobre una superfície protegida, talla la corda sota la supervisió del professor.
    \item Observa què passa.
\end{enumerate}

\textbf{Quina força ha desaparegut?} \rule{2cm}{0.4pt}

La gravetat no ha canviat. Ha desaparegut la tensió que mantenia l'objecte en equilibri.

Una corda pot estirar, però no pot empènyer. I quan la tallem, deixa de transmetre la tensió entre els seus extrems. Una manera força expeditiva de comprovar-ho.

\subsection{6. El dinamòmetre inquiet}

Torna a penjar un objecte del dinamòmetre.

\begin{enumerate}
    \item Mantén-lo quiet i anota la lectura.
    \item Eleva'l bruscament i observa l'indicador.
    \item Fes-lo baixar bruscament i torna a observar.
    \item Repeteix-ho amb moviments suaus.
\end{enumerate}

\textbf{Anota:} Quan marca més? Quan marca menys? L'indicador oscil·la?

La massa no ha canviat i la gravetat tampoc. Però la lectura pot variar quan accelerem l'aparell.

\subsection{7. Donem-hi una volta}

Amb un objecte lleuger ben subjectat al dinamòmetre, i sota la supervisió del professor, fem-lo girar amb cura descrivint una circumferència.

\textbf{Observa:} Què passa amb la direcció del moviment? El dinamòmetre continua detectant una força?

Un objecte que gira canvia contínuament la direcció de la seva velocitat, encara que la rapidesa sigui constant.

Per tant, \textbf{canviar de direcció també és canviar la velocitat}. I això requereix una força resultant.

\subsection{8. El nostre catàleg de forces}

Hem descobert unes quantes coses. Algunes eren previsibles; d'altres no tant.

Completa aquesta taula amb les observacions de la pràctica.

{\footnotesize
\setlength{\tabcolsep}{3pt}
\renewcommand{\arraystretch}{1.3}
\begin{center}
\begin{tabular}{|p{0.28\columnwidth}|p{0.58\columnwidth}|}
\hline
\textbf{Concepte} & \textbf{Què hem après?} \\ \hline
Dinamòmetre & Mesura forces en newtons. \\ \hline
Pes & Força gravitatòria sobre una massa. \\ \hline
Massa & No és el mateix que el pes. \\ \hline
Equilibri & Les forces es poden compensar. \\ \hline
Tensió & Una corda tensa transmet forces. \\ \hline
Fregament & S'oposa al lliscament. \\ \hline
Direcció & Una força té direcció i sentit. \\ \hline
Resultant & Determina el canvi de velocitat. \\ \hline
Acceleració & Canvia la velocitat o la direcció. \\ \hline
\end{tabular}
\end{center}
}

\subsection*{I ara què?}

Hem trobat forces que estiren, frenen, sostenen objectes o els fan canviar de direcció. Algunes actuen sense que res es mogui. D'altres semblen comportar-se de manera força estranya.

Durant aquest curs haurem d'explicar tot això i descobrir les lleis que governen les forces.

Newton va dedicar-hi una bona part de la seva vida. Nosaltres tenim unes quantes classes. Caldrà organitzar-se.
